\begin{abstract}
Blocking a malicious agent Skill can protect resources while leaving the user's task unfinished. Recovery becomes harder when no prepared replacement exists and some actions have already happened. We study bounded continuation construction under fixed permissions. A residual effect contract combines unfinished task goals, reusable values, attacker-created dependencies, and confirmed or uncertain side effects. A constructor proposes a continuation; a common executable checker returns counterexamples before deployment. We evaluate six held-out synthetic task structures, varying inputs, partial execution, faults, strong controls and component ablations with GPT-6 Luna and GPT-4.1 mini. With shared transport normalization, the method safely completes 55/72 and 51/72 episodes, respectively. Ordinary Luna repair also completes 55/72, but requires 103 model requests rather than 31. Shared operand guidance improves replanning but leaves a completion gap. However, a source-free deterministic compiler completes all 20 feasible matched cases, and removing several proposed components does not reduce completion. Forty-three accepted programs fail fresh inputs. These findings support explicit state constraints and reduced model use in this finite setting, while limiting claims about novel synthesis, component necessity and real-world Skill security.
\end{abstract}
